\documentclass[10pt]{article}
\usepackage[T1]{fontenc}
\usepackage[utf8]{inputenc}
\usepackage[margin=0.72in]{geometry}
\usepackage{graphicx}
\usepackage{booktabs}
\usepackage{tabularx}
\usepackage{array}
\usepackage{xcolor}
\usepackage{hyperref}
\usepackage{float}
\usepackage[font=small,labelfont=bf]{caption}
\usepackage{enumitem}
\definecolor{navy}{HTML}{183B5B}
\definecolor{pale}{HTML}{EDF4F8}
\definecolor{warn}{HTML}{FFF4E5}
\hypersetup{colorlinks=true,linkcolor=navy,citecolor=navy,urlcolor=navy}
\setlength{\parindent}{0pt}
\setlength{\parskip}{0.45em}
\setlist[itemize]{nosep,leftmargin=1.4em}
\renewcommand{\arraystretch}{1.12}
\newcommand{\evidencebox}[2]{%
  \begin{center}\fcolorbox{navy}{pale}{\begin{minipage}{0.94\linewidth}
  \textbf{#1}\quad #2
  \end{minipage}}\end{center}}
\newcommand{\warningbox}[2]{%
  \begin{center}\fcolorbox{orange!75!black}{warn}{\begin{minipage}{0.94\linewidth}
  \textbf{#1}\quad #2
  \end{minipage}}\end{center}}
\newcommand{\metricdef}{Expert beats were paired one-to-one with detections inside $\pm75$ ms. TP is a paired detection, FP is an unpaired detector output, FN is an unpaired expert beat, and $F_1=2TP/(2TP+FP+FN)$.}
\newcommand{\provenance}{The ``historical'' rows reproduce the earlier NeuroKit/Pan--Tompkins hard-gate experiment. The ``current'' rows come from the complete all-48-record audit. Small NeuroKit count differences between audit versions reflect implementation/configuration differences and must not be interpreted as physiology.}

\title{MIT--BIH Record 108\\R-Peak and RR-Reliability Case Analysis}
\author{ECG-only seizure detector: RR--HRV branch audit}
\date{4 August 2026}

\begin{document}
\maketitle

\evidencebox{Bottom line.}{The old forward-only Pan--Tompkins hard gate failed on record 108: it rejected many correct NeuroKit peaks while retaining most NeuroKit false positives. The current UNSW primary detector reduced the full-record error burden to 7 FP and 6 FN, but 20.6\% of its RR intervals still lacked the prescribed two-detector reliability support.}

\section{What record 108 contains}
The MIT--BIH Arrhythmia Database contains approximately 30-minute, two-channel ambulatory ECG records digitized at 360 samples/s with expert beat annotations \cite{moody2001mitbih,physionetmitdb,goldberger2000physiobank}. Record 108 is from an 87-year-old woman. Its channels are MLII and V1. The database documentation lists borderline first-degree AV block, sinus arrhythmia, multiform premature ventricular complexes, and considerable noise and baseline shift in the lower V1 channel \cite{physionetmitdb}.

This audit used the \emph{upper MLII channel}. Therefore the database's warning about the lower V1 channel cannot, by itself, explain an error measured in MLII. Calling every disagreement ``artifact'' would be unsupported.

\begin{tabularx}{\linewidth}{@{}>{\bfseries}p{0.22\linewidth}X@{}}
Sampling & 360 Hz; one sample is 2.78 ms.\\
Reference & MIT--BIH \texttt{atr} beat annotations; symbols identify beat classes, not detector correctness.\\
Primary question & Does a second QRS detector safely validate NeuroKit detections, and does the replacement primary detector improve the result?\\
Interpretation limit & One annotated database record is a detector case study, not evidence of seizure detection or patient-generalized performance.\\
\end{tabularx}

\section{Full-record detector results}
\metricdef\ \provenance

\begin{table}[H]
\centering
\caption{Record 108 full-record results.}
\begin{tabular}{@{}llrrrr@{}}
\toprule
Audit & Method & TP & FP & FN & $F_1$\\
\midrule
Historical & NeuroKit alone & 1,563 & 181 & 200 & 0.891\\
Historical & NeuroKit + forward PT veto & 758 & 161 & 1,005 & 0.565\\
Current & NeuroKit 300-ms baseline & 1,564 & 181 & 199 & 0.8917\\
Current & UNSW initial event & 1,757 & 7 & 6 & 0.9963\\
Current & UNSW refined R fiducial & 1,757 & 7 & 6 & 0.9963\\
\bottomrule
\end{tabular}
\end{table}

\begin{figure}[H]
\centering
\includegraphics[width=0.82\linewidth]{figures/record_108_current_metric_comparison.png}
\caption{Current full-record comparison. Error counts, rather than $F_1$ alone, show the type of failure.}
\end{figure}

\section{Why the old hard gate failed}
The tested rule retained a NeuroKit peak only when a Pan--Tompkins event occurred 0--100 ms \emph{after} it. Pan--Tompkins detects QRS energy through filtering, differentiation, squaring, moving-window integration, and adaptive thresholds \cite{pan1985qrs}; it is not an expert label and its event time need not equal the R-wave apex. Here the direction and width of the gate were mismatched to many true beats.

The gate retained only $758/1{,}563=48.5\%$ of NeuroKit's true detections but retained $161/181=89.0\%$ of its false detections. It therefore selected the wrong subset: FN increased from 200 to 1,005 while FP fell by only 20.

\begin{figure}[H]
\centering
\includegraphics[width=\linewidth]{figures/record_108_historical_hard_gate.png}
\caption{Historical hard-gate experiment. The same strip is plotted twice. Orange triangles in the right panel are expert beats lost by the 0--100-ms forward veto.}
\end{figure}

Published two-detector RR-quality work has used symmetric support around a primary event rather than this forward-only veto \cite{ho2024rrquality,ho2025rraccuracy}. That supports testing bidirectional timing, but it does not prove that a wider gate is correct: broad tolerances can hide clinically relevant timestamp error \cite{porr2024jitter}. Detector agreement is context, not ground truth.

\section{What the complete architecture changed}
The current branch uses the Khamis/UNSW detector as the primary QRS-event source \cite{khamis2016telehealth,kristof2024qrs}. NeuroKit's gradient detector remains a secondary context stream, and the initial event is refined to the highest positive amplitude within $\pm50$ ms. A primary event is ``supported'' only when exactly one secondary event is inside $\pm50$ ms. An RR interval is supported only when both boundary events are supported and exactly two secondary events occur across the expanded RR span, following the published two-detector structure \cite{ho2024rrquality,ho2025rraccuracy}.

\begin{figure}[H]
\centering
\includegraphics[width=\linewidth]{figures/record_108_complete_architecture_strip.png}
\caption{Current architecture on an explanatory strip. Diamonds are expert beats; triangles and crosses are detector outputs. Their vertical separation can reflect fiducial conventions in a biphasic complex; matching is performed in time, not amplitude.}
\end{figure}

Across the full record, 1,400 of 1,763 output RR intervals were supported (coverage 79.41\%). The mean absolute RR error among all evaluable intervals was 17.46 ms; among supported evaluable intervals it was 12.33 ms. Support enriched for lower RR error, but did not make the timestamps exact.

\warningbox{Remaining concern.}{The current detector is excellent on this annotated record but is not yet validated on peri-ictal ECG. Before HRV features are trusted, representative patient ECG must be manually annotated and performance stratified by rhythm, morphology, noise, and heart rate.}

\section{Decision for the RR--HRV branch}
\begin{itemize}
\item Reject the 0--100-ms forward Pan--Tompkins hard veto.
\item Retain UNSW initial QRS events as the current primary candidate for this record.
\item Retain detector support and RR coverage as numerical context; do not label unsupported intervals ``artifact'' automatically.
\item Keep the R-fiducial refinement under audit, even though it did not change record-level counts here.
\end{itemize}

\begingroup\interlinepenalty=10000
\bibliographystyle{unsrt}
\bibliography{MITBIH_Record_Case_Reports}
\endgroup
\end{document}
